% SPDX-License-Identifier: CC-BY-4.0
% Copyright (c) 2026 Martin Schr\"oder <info@swedishembedded.com>

\section{Reproducibility}\label{sec:reproducibility}

The software is the open-source Splinter pipeline; the study is its sample
\code{samples/synthetic-author}. An author, its corpus plan, its examination and every
reference answer are functions of the seed and the planned article count; nothing generated
is checked in. Model references are given as placeholders: \code{<hosted-model>} is the small
hosted language model that writes the essays, reads them for the gate and extracts the
decisions (reached with \code{-{}-allow-remote}); the policy is Qwen3-8B at \code{\$QWEN3\_8B}.
The hosted model is sampled at its provider's defaults, so a rerun writes different essays;
the plan, the oracle, the examination and the grading are reproduced exactly.

\subsection{Subcommands}

\begin{description}
\item[\code{persona}] draws the author of a seed and prints its name, its beliefs in order and a few of its decisions.
\item[\code{plan}] lists the cases each article of a corpus tells, with the decision, deciding
  belief, facts and options of each.
\item[\code{corpus}] writes the essays under the gates of \cref{sec:corpus}, within a budget of
  requests and estimated output tokens.
\item[\code{mine}] reads the decisions out of the essays (\cref{sec:mining}).
\item[\code{dataset}] makes the training records (\cref{sec:records}).
\item[\code{train}] holds out a tenth of the cases and trains the adapter (\cref{sec:lora}).
\item[\code{exam}] puts the examination to one arm and appends each answer to
  \code{exam/<arm>.jsonl}; the first answer recorded for an item stands.
\item[\code{report}] tabulates the arms of one author with Wilson intervals and the six paired
  tests.
\item[\code{cross}] prints the divergence table of \cref{tab:cross}.
\end{description}

\subsection{Commands}

For seed $S\in\{101,202,303\}$, with a corpus directory \code{\$C} and an examination
directory \code{\$E} whose \code{articles/} is a copy of the corpus's \code{articles/} (the
retrieval arm indexes it):
\begin{verbatim}
synthetic-author persona --seed $S
synthetic-author corpus --seed $S --cache $C --writer <hosted-model> \
  --reader <hosted-model> --allow-remote --articles 60 \
  --episodes-per-article 3 --min-words 700 --attempts 4 \
  --max-calls <N> --max-output-tokens <M> --workers <W>
synthetic-author mine --cache $C --miner <hosted-model> --allow-remote
synthetic-author dataset --seed $S --cache $C
synthetic-author plan --seed $S --articles 60 > plan.jsonl
python3 scripts/audit.py plan.jsonl $C/decisions.jsonl
synthetic-author train --cache $C --base $QWEN3_8B --steps 240 \
  --rank 16 --alpha 32 --records-per-step 4 --bf16
for ARM in base named principles retrieval; do
  synthetic-author exam --seed $S --cache $E --base $QWEN3_8B --arm $ARM \
    --articles 100 --per-cell 40 --max-tokens 256
done
synthetic-author exam --seed $S --cache $E --base $QWEN3_8B --arm tuned \
  --adapter $C/train/adapter.safetensors --articles 100 --per-cell 40 \
  --max-tokens 256
synthetic-author report --cache $E
synthetic-author cross --exams <dir> --seeds 101,202,303 --articles 100
\end{verbatim}
The corpora were planned with 60 articles. The examination and the divergence table are drawn
with \code{-{}-articles 100}: that value fixes which case names an item may not take, so it must
be the same for \code{exam} and \code{cross}, and it is the value the reported examinations
were drawn with. Here \code{<dir>} holds one examination directory per author, named
\code{author-<seed>}. The audit matches each decision to the plan by article and case with
\code{scripts/audit.py}.

\subsection{Data and tables}

\begin{description}
\item[\code{data/arms.csv}, \code{data/cells.csv}, \code{data/paired.csv}] items right per arm,
  per cell and per pair, made from the examination outputs by
  \code{scripts/collect.py EXAMS\_DIR data 101,202,303}, where \code{EXAMS\_DIR} holds
  \code{author-<seed>/exam/<arm>.jsonl}.
\item[\code{data/cross.csv}] the table printed by \code{cross}.
\item[\code{data/audit.csv}] the counts printed by \code{scripts/audit.py}.
\item[\code{data/corpora.csv}] per corpus: planned articles, essays kept and given up, words,
  writing and reading requests (whose sum is the request count in the corpus directory's
  \code{usage.json}), the lines of \code{rejected.jsonl} by reason, and the essays mined,
  cases read, decisions kept, fixed-wording briefings and records of \code{mine} and
  \code{dataset}.
\item[\code{data/training.csv}] the records trained on and held out (\code{train/train.jsonl}
  and \code{train/held\_out.jsonl}), the loss and elapsed time logged at step 240, the
  learning rate and the held-out scores of the training record where it was written.
\item[\code{data/authors.csv}] the output of \code{persona}.
\item[\code{data/cost.csv}] the estimated token counts of \code{usage.json}.
\item[\code{data/spend.csv}] the provider's reported daily spend of the account before and
  after the corpus runs.
\end{description}
\code{scripts/build\_data.py} (standard library only) computes the Wilson intervals, the exact
sign tests and the Newcombe intervals from these files. It refuses to continue unless the
tables are consistent with one another (cells add up to arms, discordant counts agree with
the right counts, the request and rejection counts of each corpus reconcile, records equal
three per briefing, the audit covers the decisions kept, trained and held-out records add up
to the records made) and the derived values equal those recorded when the examinations were
run. It writes the LaTeX tables, the macros and the plot data under the build directory.
\code{make} in \code{papers/persona-training} builds \code{target/papers/persona-training.pdf};
nothing is written into the paper's directory.
